% Cross-Source Calibration Methodology — Thematic Synthesis
% !TEX program = xelatex
% !TEX encoding = UTF-8 Unicode

\section{Cross-Source Calibration Methodology}
\label{sec:calibration}

One of the key findings of this project is that data from one source can
resolve ambiguity in another. This cross-source calibration is the most
powerful tool in the conversion toolkit. Because each recorder had different
strengths and weaknesses --- Edwards was reliable for /h/ but not for vowel
quality, Williams was reliable for stress but not for preaspiration, Eliot
was systematic but English-filtered --- the sources complement each other.
By calibrating across sources, we can recover phonological information that
no single source provides.

\subsection{The Calibration Pipeline}

The calibration pipeline proceeds from most reliable to least reliable, with
each source used to calibrate the next. The ordering is determined by the
recorder's linguistic training, fluency, and orthographic system:

\begin{enumerate}
  \item \textbf{Fielding (2013)} --- Modern phonemic orthography with 1:1
    phoneme-to-grapheme correspondence. The most phonologically reliable
    source, but limited to Mohegan-Pequot and representing the modern
    revived language.

  \item \textbf{Trumbull/Eliot (1903/1666)} --- Systematic missionary
    orthography with the largest corpus (4,491 records). Eliot's grammar
    provides explicit phonological descriptions, and the large corpus enables
    statistical approaches.

  \item \textbf{Edwards (1787)} --- Fluent L2 speaker with the most reliable
    /h/ detection. Small corpus (144 records) but high per-record
    reliability.

  \item \textbf{Prince-Speck (1904)} --- Trained linguist with explicit /hC/
    cluster notation. Medium corpus (446 records) but two-stage recording
    chain introduces compounding errors.

  \item \textbf{Williams (1643)} --- Linguistically gifted amateur with
    unique stress diacritics. Large corpus (2,134 records) but high vowel
    ambiguity and no /h/ detection.

  \item \textbf{Wood/Winslow (1634/1624)} --- Raw English perception with no
    linguistic training. Small corpora (339 and 24 records) but valuable as
    the earliest attestations.
\end{enumerate}

\subsection{Cross-Source Disambiguation: Specific Examples}

\subsubsection{Massachusett Disambiguates Narragansett Vowels}

Williams' vowel ambiguity in Narragansett is partially resolved by comparing
against Eliot's Massachusett cognates, which have a more systematic
orthography. For example, Williams' \textit{a} can represent /ə/, /ɔː/, or
/aː/. When the Massachusett cognate shows \textit{â} (Eliot's long a marker),
the Narragansett vowel is likely /aː/. When the Massachusett cognate shows
\textit{u} (Eliot's schwa marker), the Narragansett vowel is likely /ə/.

The cognate pair for 'dog' illustrates this: Williams writes \textit{ayim},
Eliot writes \textit{arum}. The PA reconstruction *aɬemwa confirms the
vowel quality and shows the /ɬ/ $\rightarrow$ /j/ (Narragansett) vs.\ /ɬ/
$\rightarrow$ /r/ (Massachusett) reflex difference. The cross-source
comparison resolves both the vowel quality (first vowel /a/) and the
consonant reflex.

\subsubsection{Edwards Disambiguates Preaspiration}

Edwards' consistent /h/ marking in Mahican provides a model for predicting
preaspiration in Williams and Wood, where it is unwritten. When Edwards
writes \textit{hkeesque} 'eye' with initial /hk/, and the PA reconstruction
shows */h/ in a corresponding position, we can infer that the cognate
Narragansett form (which Williams writes without /h/) also had preaspiration.
This inference is strengthened when the morphotactic position matches: a stop
in the coda of a stressed syllable is a candidate for preaspiration.

\subsubsection{Prince-Speck Confirms Fielding}

Where Prince-Speck (1904) and Fielding (2013) agree on a Mohegan-Pequot
form, the confidence is high. For example, both record \textit{wetu} for
'house', showing remarkable lexical stability across 110 years. Where they
disagree, the circularity risk must be assessed: Fielding's dictionary is
partially reconstructive, and some entries may derive from Prince-Speck
itself. Using Fielding to ``verify'' Prince-Speck for such entries would be
circular. The conversion workflow must track which Fielding entries derive
from Prince-Speck (self-source) vs.\ from independent sources.

\subsubsection{Wood and Winslow Corroborate Each Other}

Wood (339 records) and Winslow (24 records) are the two earliest Wampanoag
sources, separated by only 10 years. Where they agree on a form, the
reading is likely correct. For example, both record \textit{matta} for 'no',
and both record \textit{sachim} for 'chief'. Where they disagree --- as with
Wood's \textit{abamocho} vs.\ Winslow's \textit{hobbamock} for 'devil' ---
the disagreement may reflect real dialect differences (Wood records
Massachusetts Bay, Winslow records Pokanoket) or different lexical roots
entirely.

\subsection{The Cognate Set Approach}

The most powerful calibration tool is the cognate set: a word with the same
meaning attested in multiple SNEA languages, with PA reconstruction as the
anchor. Table~\ref{tab:cognate-set} shows a representative cognate set for
basic vocabulary items.

\begin{table}[h]
\centering
\caption{Cross-source cognate set for basic vocabulary}
\label{tab:cognate-set}
\small
\begin{tabular}{llllll}
\toprule
\textbf{Gloss} & \textbf{Williams} & \textbf{Eliot} & \textbf{Edwards} & \textbf{Fielding} & \textbf{PA} \\
 & \textbf{(NRG)} & \textbf{(WMP)} & \textbf{(MAH)} & \textbf{(MPQ)} & \\
\midrule
'house' & wétu & wetu & --- & wetu & *wiːkiwaːmi \\
'water' & nipp & nippe & --- & nip & *nepyi \\
'dog' & ayim & arum & --- & ayum & *aɬemwa \\
'man' & weenawash & wénawaš & --- & --- & --- \\
'fire' & apeisuwau & apeisuwâ & --- & --- & --- \\
'I' & neen & neen & neen & neen & *niːra \\
'two' & nees & nees & neesoh & nees & *niːšwi \\
'no' & matta & matta & --- & matta & *mati- \\
'woman' & squa & squa & --- & squa & *eːθkweːwa \\
'five' & --- & pàpàna & nannah & --- & *nyaːlanwi \\
\bottomrule
\end{tabular}
\end{table}

The cognate set reveals several patterns:

\begin{enumerate}
  \item \textbf{Lexical stability:} Basic vocabulary items like 'house',
    'water', 'I', and 'two' are nearly identical across all SNEA languages
    and across 260+ years (Williams 1643 to Fielding 2013). This confirms
    the reliability of the core vocabulary for cross-source calibration.

  \item \textbf{Regular sound correspondences:} The PA *r reflex distinguishes
    N-dialects (Massachusett \textit{neen} from *niːra) from Y-dialects
    (Mohegan \textit{neen} with /n/ from a different PA source). The PA *ɬ
    reflex distinguishes Narragansett /j/ (\textit{ayim}) from Massachusett
    /r/ (\textit{arum}).

  \item \textbf{Cross-source vowel confirmation:} The vowel in 'house' is
    consistently \textit{e} across Williams, Eliot, and Fielding, confirming
    the /i/ quality (from PA *iː). The vowel in 'water' is consistently
    \textit{i} across all sources, confirming the /ɪ/ quality (from PA *e).

  \item \textbf{Edwards' Mahican as an outlier:} Edwards' \textit{neesoh}
    'two' shows the Mahican final vowel (\textit{-oh}) that is absent in the
    other SNEA languages, reflecting Mahican's transitional status between
    SNEA and Delaware.
\end{enumerate}

\subsection{How PA Reconstructions Anchor Vowel Quality}

Proto-Algonquian reconstructions (Bloomfield 1925, 1946; Aubin 1975;
Pentland 1979) provide the comparative evidence needed to choose among
ambiguous vowel mappings. The logic is:

\begin{enumerate}
  \item For a given colonial spelling, identify the set of possible phonemic
    interpretations (from the grapheme-to-phoneme mapping table).

  \item Find the PA reconstruction for the cognate set.

  \item Apply the known sound correspondences (PA $\rightarrow$ daughter
    language) to predict the expected reflex.

  \item Select the phonemic interpretation that matches the expected reflex.
\end{enumerate}

For example, Williams' \textit{wétu} 'house' has \textit{e} in the first
syllable. The possible phonemic interpretations of Williams' \textit{e} are
/iː/, /ɛ/, /ə/, or ∅. The PA reconstruction is *wiːkiwaːmi. The known
correspondence is PA *iː $\rightarrow$ Narragansett /iː/. Therefore, the
\textit{e} in \textit{wétu} represents /iː/, not /ɛ/ or /ə/. This is
confirmed by the consistent \textit{ee} spelling in Eliot's Massachusett
cognate \textit{wetu} (where \textit{ee} = /iː/).

This logic applies to every ambiguous vowel in the corpus. The PA
reconstruction is the anchor that resolves the ambiguity.

\subsection{Calibration Confidence Levels}

Each converted form is assigned a confidence level based on the strength of
the calibration evidence:

\begin{itemize}
  \item \textbf{[high]} --- Confirmed by cognate in a more reliable source or
    by PA reconstruction. No contradictory evidence.

  \item \textbf{[med]} --- Default rule applies, no contradicting cognate,
    but no confirming cognate either.

  \item \textbf{[low]} --- Guess; no cognate available, default rule may not
    apply.

  \item \textbf{[FLAG]} --- Contradictory evidence; manual review needed.
\end{itemize}

The calibration pipeline produces a confidence-tagged output that allows
downstream users to assess the reliability of each converted form. This
transparency is essential for both scholarly analysis and revitalization
applications, where the confidence level determines how the form can be
used.
